
%%%%%%%%%%%%%%%%%%%%
\begin{frame}[fragile]{La commande \sqlkw{create table}}

Un premier exemple: la table \texttt{Voyageur}.

\vfill

\begin{minted}[baselinestretch=.9,
               fontsize=\small,
               framesep=2mm]{sql}
create table Voyageur (idVoyageur integer not null,
                      nom varchar(30) not null,
                      prénom varchar(30)  not null,
                      ville varchar(30) not null,
                      région varchar(30) not null,
                      primary key (idVoyageur)
                   )
\end{minted}

\vfill
Quelques choix à effectuer: conventions de nommage, accents, etc.

\end{frame}


%%%%%%%%%%%%%%%%%%%%
\begin{frame}[fragile]{La contrainte \sqlkw{not null}}

Les valeurs à \sqlkw{null} sont source de problèmes: on peut les interdire
avec \sqlkw{not null} ou donner une valeur par défaut.

\vfill

\begin{minted}[baselinestretch=.9,
               fontsize=\small,
               framesep=2mm]{sql}
create table Logement 
  (code varchar(4) not null,
   nom varchar(50) not null,
   capacité   integer not null,
   type varchar(10) default 'Hôtel',
   lieu varchar(30) not null,
   primary key (code)
)
\end{minted}

\vfill
Le SGBD rejettera alors toute tentative d’insérer un nuplet avec une valeur manquante.

\end{frame}

%%%%%%%%%%%%%%%%%%%%
\begin{frame}[fragile]{La clé primaire}

Elle est spécifiée avec la clause  \sqlkw{primary key}.

\vfill

\begin{minted}[baselinestretch=.9,
               fontsize=\small,
               framesep=2mm]{sql}
create table Activité 
   (codeLogement varchar (4) not null,
    codeActivité      varchar(30) not null,
    description       varchar(255) not null,
    primary key (codeLogement, codeActivité),
    foreign key (codeLogement) 
       references Logement(code)
)
\end{minted}

\vfill
Une clé primaire peut comprendre plusieurs attributs.
\vfill
Tous les attributs d'une clé primaire doivent être \sqlkw{not null}

\end{frame}


%%%%%%%%%%%%%%%%%%%%
\begin{frame}[fragile]{Clé primaire, clés ``secondaires''}

On peut définir d'autres clés avec la clause \sqlkw{unique}.

\vfill

\begin{minted}[baselinestretch=.9,
               fontsize=\small,
               framesep=2mm]{sql}
create table Logement 
   (code varchar(4) not null,
    nom varchar(50) not null,
    capacité   integer not null,
    type varchar(10) default 'Hôtel',
    adresse varchar(30) not null,
    primary key (code),
    unique (adresse)
)
\end{minted}

\vfill

Permet d'avoir un identifiant ``abstrait'', non modifiable, et d'ajouter des
contraintes flexibles sur les attributs descriptifs.

\end{frame}


%%%%%%%%%%%%%%%%%%%%
\begin{frame}[fragile]{Clé étrangère}

Elle est spécifiée avec la clause  \sqlkw{foreign key}.

\vfill
\begin{minted}[baselinestretch=.9,
               fontsize=\small,
               framesep=2mm]{sql}
create table Séjour  (idSéjour integer not null,
                      idVoyageur integer not null,
                      codeLogement varchar (4) not null,
                      début integer not null,
                      fin integer not null,
                      primary key (idSéjour),
                      foreign key (idVoyageur) 
                         references Voyageur(idVoyageur),
                      foreign key (codeLogement) 
                         references Logement(code)
                      )
\end{minted}

\vfill

Si la clé étrangère est \sqlkw{not null}: association de type \texttt{1-n},
sinon association de type \texttt{0-n}.

\end{frame}

%%%%%%%%%%%%%
\begin{frame}{À retenir}

Après la phase de conception, la création du schéma ne présente aucune difficulté.
\vfill

\begin{redbullet}
   \item  Connaître les principaux types SQL (cf. support de cours)
   \item Connaître la commande \sqlkw{create table}
   \item Ajouter les contraintes de clés primaire et étrangères.
\end{redbullet}

\vfill

La spécification des contraintes n'est pas une lourdeur. Elle garantit que la base est saine,
et évité beaucoup d'ennuis.
\end{frame}


